\subsection{Does the vote improve prediction beyond deletion size?}
\label{sec:incremental}
Two proposals removing the same number of columns can have different refit outcomes. To test whether their mask scores carry information beyond width, we predicted XGBoost refit failure on a task using models trained on the other thirteen. The first logistic model used the fraction removed and linear spline terms with knots at $1/4$, $1/2$ and $3/4$. The second added the largest fraction of dissenting panels in any separating restoration. We trained and evaluated each collection separately, weighting tasks equally, then folds within task and proposals within fold equally. All scaling and coefficient fitting used the training tasks with a fixed ridge penalty of 0.01 and an unpenalized intercept. The primary measure was held-out binary log loss, including folds in which every refit had the same outcome.

Adding dissent changed mean log loss from 0.632 to 0.634 on paths (paired change $+0.002$, 95\% interval $[-0.030,0.049]$) and from 0.580 to 0.567 on drawn proposals ($-0.013$, $[-0.034,0.006]$). The evidence for better probability prediction was inconclusive. Within the 28 folds per collection containing both outcomes, candidate-ranking AUC rose from 0.784 to 0.808 on paths and from 0.779 to 0.825 on drawn deletions. The paired increases were 0.024 ($[0.006,0.045]$) and 0.046 ($[0.002,0.102]$). At equal widths the size model always ties. The augmented model ranked failing refits above passing ones with concordance 0.582 on paths (97 pairs in 22 folds) and 0.573 on drawn deletions (137 pairs in 27 folds). The increases over 0.5 had intervals $[0.024,0.144]$ and $[-0.001,0.150]$, respectively. We used half credit for ties and 50,000 paired task resamples, holding out-of-task predictions fixed.

\paragraph{Converting predicted risks to vetoes.}
At unit false-alarm cost and miss cost $r$, we vetoed when the held-out predicted failure probability reached $1/(1+r)$. This threshold equates the costs under the estimated probability. With no further threshold fitting, the total costs for size alone and size plus votes at $r=1,2,5,10$ were, respectively, 155/148, 190/210, 214/222 and 211/223 on paths. They were 130/132, 156/134, 184/178 and 216/199 on drawn deletions. At $r=2$ the vote reduced drawn cost on six held-out tasks and tied on eight. Removing each task in turn and refitting the remaining leave-one-task-out models gave reductions of 1.14--6.41 cost units per 100 proposals. These gains are relative to the fitted size model. At $r=1$ the simple size cutoff chosen on other tasks made 98 drawn errors, compared with 132 for the augmented model.

\subsection{Which candidate should receive a limited refit?}
\label{sec:budgeted}
An ordering of proposals has practical value when only some can be refitted. We split the ranking rows equally into proposal and guard samples, stratified by class. One Wine Quality fold had a class represented by one ranking row, leaving 41 feasible folds on all 14 tasks. Within a fold all methods shared the XGBoost fit, the anchor determined on proposal rows, and a pool of at most eight candidates per width. At each distinct strictly smaller width among $\lceil |A|/4\rceil$, $\lceil |A|/2\rceil$ and $\lceil3|A|/4\rceil$, the pool included the Gain prefix and up to seven distinct random subsets of $A$.

The primary budget allowed one candidate refit at half width, rounded up. Gain nominated the highest-ranked prefix. Voting nominated the candidate with the most approvals in its worst restoration, breaking ties by the smallest sum of Gain ranks. This nomination can choose a set that fewer than three panels approved. At the secondary budget each rule nominated one set per available width. We also used a common random tie order to measure the effect without Gain. Candidate and anchor models fitted on fitting and proposal rows predicted on the separate guard sample. We kept the smallest nominee passing the guard, or $A$ if all failed, and refitted the final model on all development rows. The candidate-refit budget was identical within folds. One fold had a one-column anchor and offered no deletion.

We measured the fraction of anchor columns saved when the final refit met its outer tolerance, $U=(1-|S|/|A|)\mathbf 1\{\text{outer refit passes}\}$. A fallback and a deletion that breaches tolerance both receive zero. Table~\ref{tab:operating-guards} also reports breaches separately so that accepting a harmful deletion can carry its own cost.

\begin{table}[tb]
\centering
\caption{Comparing proposals at the same guard-refit budget. $U$ is the fraction of anchor columns saved when the final refit passes the outer tolerance, and zero otherwise. Breaches count folds; other columns weight tasks equally.}
\label{tab:operating-guards}
\small\setlength{\tabcolsep}{3pt}
\begin{tabular}{@{}llrrr@{}}
\toprule
Budget & Within-width order & $100\bar U$ & Breaches & AUC \\
\midrule
One & Gain & 9.33 & 1/41 & 0.9351 \\
 & Vote, then Gain & 10.12 & 3/41 & 0.9341 \\
 & Lottery & 11.29 & 5/41 & 0.9307 \\
 & Vote, then lottery & 12.48 & 4/41 & 0.9321 \\
\midrule
Up to three & Gain & 15.36 & 2/41 & 0.9305 \\
 & Vote, then Gain & 16.75 & 6/41 & 0.9279 \\
 & Lottery & 11.55 & 10/41 & 0.9286 \\
 & Vote, then lottery & 12.66 & 11/41 & 0.9281 \\
\bottomrule
\end{tabular}
\end{table}


At half width all candidates scored zero approvals in their worst restoration on 32 of the 40 folds offering a deletion. Voting changed five nominations and four final supports across 41 folds. Mean $U$ rose from 9.33\% to 10.12\% of the anchor (paired gain 0.79 percentage points, whole-task bootstrap 95\% interval $[0,2.38]$). One New Thyroid fold accounted for the entire increase. The two-column Gain prefix received zero of four approvals in its worst restoration and failed the guard. A different pair received four of four and passed. Its outer AUC was 0.9982 and accuracy 97.22\%, versus 0.9943 and 95.83\% for the three-column anchor.

On SPECTF the best seven-column nominee received only one of four approvals in its worst restoration. The guard nevertheless admitted it, and outer AUC fell from 0.9662 for the 14-column anchor to 0.9222. On Ionosphere a second additional breach reduced accuracy from 94.02\% at 14 columns to 90.60\% at seven. Across all folds, breaches rose from one to three and mean AUC fell from 0.9351 to 0.9341. Keeping only outer rows whose predictors were absent from development preserved all 41 pass/fail decisions and $U$ values for these two policies. This sensitivity changes scoring conditional on the original selections.

We timed the two half-width policies separately on one logical CPU, with fresh fits and caches in two order-balanced pairs per fold. Each clock ran from in-memory preprocessing through final prediction, including nomination and the guard decision. The prefix constructed only its own candidate. The vote stopped evaluating a candidate once its worst approval count could no longer beat the incumbent under the Gain tie rule. Across 164 runs the task-weighted mean time was 0.8423 seconds for Gain and 0.8658 for voting, a difference of 0.0235 seconds (whole-task 95\% interval $[-0.0584,0.1046]$). Voting averaged 11.73 additional mask evaluations. The nominees, final supports and predictions agreed with the full evaluations.

For these 14 tasks, let $b$ value one unit of $U$, let $h$ price an outer breach and let $c$ be cost per second. Using the measured point estimates, the half-width vote has positive net benefit over Gain when $0.00794 b>0.04762 h+0.02346 c$. With computation free this requires $h/b<1/6$. For example, pricing a breach at $b/12$ and setting $c=0$ gives a net gain of 0.40 points per unit $b$ (whole-task 95\% interval $[-0.79,2.18]$ points). Omitting New Thyroid changes the point estimate to $-0.43$ points. With up to three candidate refits, mean $U$ increased by 1.39 points ($[-1.50,4.76]$), breaches rose from two to six and mean accuracy fell from 82.02\% to 81.45\%.

\paragraph{A held-out handwritten-digit test.}
We fixed the width rules before scoring UCI Optical Recognition of Handwritten Digits (Optdigits)~\cite{optdigitsdata}. The official 3,823 training and 1,797 test examples came from 30 and 13 distinct writers, respectively. We used 75\% of training rows for fitting and split the remainder equally between proposal and guard samples. Supports and test predictions were saved before the scorer read test labels. Results weight rows equally because the feature files omit writer IDs. All 24 candidates had zero approvals in their worst restoration. The one-refit policies therefore nominated the same Gain prefix. Its rejection left the 24-column anchor (AUC 0.99785, accuracy 94.66\%). At the three-refit budget both policies kept 18 columns within tolerance (0.99678, 92.65\%). The complete vote evaluated nine additional masks at the first budget and 25 at the second.

\paragraph{Learning which width should receive one refit.}
After the fixed-width comparisons we tested one exploratory rule across the entire candidate pool. For each candidate let $d=1-|S|/|A|$ and let $\hat p$ be its predicted breach probability from Sect.~\ref{sec:incremental}, trained on other tasks. We nominated the set maximizing $d(1-\hat p)$, breaking ties by Gain, then applied the same one-candidate guard. This score credits a larger deletion if predicted breach risk is low. It omits the probability of passing the guard, which also determines realized $U$. We compared the size-only and vote-augmented models on the same candidates. Models trained on drawn deletions saw $d\le1/2$ and thus extrapolated for some quarter-width candidates. Every selected nominee lay within the training range in this comparison.

In the primary comparison, the models learned from drawn deletions. Across 41 folds, adding dissent changed 16 nominations (including three changes of width) and 15 final supports. Seven folds gained compression within tolerance and three lost it. Mean $U$ rose from 8.94\% to 11.92\% (paired gain 2.98 points, whole-task 95\% interval $[-1.47,7.46]$), breaches rose from one to four, and AUC fell from 0.93291 to 0.93042. Training on path proposals instead gave mean $U$ of 7.64\% and 10.28\% (gain 2.64 points, $[-1.65,6.94]$) with one and three breaches. Both models and training sources chose the same 18-column support when applied retrospectively to the previously opened Optdigits test.
